\chapter{Úvod}
\label{sec:Uvod}

% Zdroje použité pre túto kapitolu:
% - writing_guide.md (štruktúra)
% - updates_structured.md Fáza 1 (motivácia)
% - thesis_structure.md (obsah)
% - project_analysis.md (technický prehľad)
% - zameranie_praxe.md (pôvodná pozícia DevOps)

%% MOTIVÁCIA PRE VÝBER PRAXE %%

Pôvodne som sa na odbornú prax vo firme Railsformers~s.r.o. hlásil na pozíciu DevOps špecialistu. Náplňou mala byť migrácia projektu do verejného cloudu Azure, príprava Docker kontajnerov a~vylepšovanie CI/CD pipeline v~GitLabe. Po úvodných konzultáciách s~mentorom sa však naskytla príležitosť pracovať na projekte, ktorý prirodzene kombinoval DevOps aspekty (Docker Compose orchestrácia, kontajnerizácia služieb, správa infraštruktúry) s~full-stack vývojom (Python backend, React frontend, návrh databázy). Táto kombinácia ma zaujala viac než čisto infraštruktúrna práca, a~preto som ponuku prijal.

Projekt zameraný na IoT sledovací systém pokrýval celé spektrum moderného softvérového vývoja -- od práce s~hardvérom a~IoT protokolmi, cez návrh databázovej architektúry a~backend development, až po tvorbu interaktívneho webového rozhrania. Ďalšou významnou motiváciou bola príležitosť pracovať s~technológiami, ktoré nemajú široké zastúpenie v~bežnej výuke, ale sú v~praxi stále častejšie využívané. Technológia LoRaWAN pre bezdrôtovú komunikáciu na dlhé vzdialenosti, protokol MQTT pre IoT messaging alebo špecializované databázové rozšírenie TimescaleDB pre časové rady predstavovali pre mňa atraktívnu výzvu na rozšírenie znalostí.

Firma Railsformers má s~týmito technológiami praktické skúsenosti -- ich infraštruktúra Hedurio, na ktorej beží okrem iného MQTT broker použitý v~tomto projekte, slúži ako základ pre rôzne IoT aplikácie. Vďaka tomu som očakával kvalitný mentoring, čo sa aj potvrdilo. Nemenej dôležitým faktorom bola možnosť pracovať na projekte s~reálnou komerčnou hodnotou. Sledovací systém pre eventy má praktické využitie v~oblasti bezpečnosti, logistiky a~organizácie akcií. Možnosť vidieť, ako teoretické znalosti zo štúdia prechádzajú do funkčného produktu, bola pre mňa silnou motiváciou pre výber práve tejto praxe.

%% KONTEXT A PROBLEMATIKA IoT %%

Internet vecí (IoT -- Internet of Things) predstavuje jedno z~najdynamickejšie rastúcich odvetví informačných technológií. Podľa odhadov spoločnosti Statista dosiahne počet pripojených IoT zariadení do roku 2030 viac než 29~miliárd kusov celosvetovo~\cite{statista_iot}. Tieto zariadenia generujú obrovské množstvo dát, ktoré vyžadujú špecifické prístupy k~ukladaniu, spracovaniu a~vizualizácii.

V~kontexte sledovacích systémov pre osoby hrá kľúčovú úlohu technológia LoRaWAN (Long Range Wide Area Network). Na rozdiel od Wi-Fi alebo Bluetooth, ktoré majú dosah v~ráde desiatok metrov, LoRaWAN umožňuje komunikáciu až na vzdialenosť približne 15~kilometrov v~optimálnych podmienkach pri zachovaní nízkej spotreby energie~\cite{lorawan_spec}. To ju robí ideálnou pre aplikácie, kde je nutné pokryť rozsiahle oblasti ako festivalové areály, skladové komplexy alebo mestské štvrte, bez potreby budovať hustú sieť prístupových bodov.

Existujúce komerčné riešenia pre GPS tracking sú často nákladné, uzavreté a~neposkytujú dostatočnú flexibilitu pre špecifické požiadavky zákazníkov. Open-source alternatívy zas vyžadujú značné úsilie pri integrácii jednotlivých komponentov. Projekt, na ktorom som pracoval, mal za cieľ vytvoriť modulárne riešenie kombinujúce výhody oboch prístupov -- vlastný vývoj s~využitím otvorených technológií pri zachovaní profesionálnej kvality a~rozšíriteľnosti.

%% CIELE PRÁCE %%

Hlavným cieľom práce bolo navrhnúť a~implementovať GPS tracking systém postavený na technológii LoRaWAN, ktorý umožní real-time sledovanie osôb vybavených GPS trackermi. Systém mal spracovávať dáta z~trackerov SenseCap T1000-B komunikujúcich cez LoRaWAN gateway, ukladať ich do databázy optimalizovanej pre časové rady a~zobrazovať ich v~interaktívnej webovej aplikácii.

Medzi čiastkové ciele patrilo zabezpečenie spoľahlivého príjmu a~parsovania dát z~trackerov, návrh škálovateľnej databázovej schémy s~využitím rozšírení PostgreSQL (TimescaleDB pre časové rady, PostGIS pre priestorové dotazy), implementácia REST~API pre historické dáta aj WebSocket komunikácie pre real-time aktualizácie, a~vytvorenie používateľsky prívetivého webového rozhrania s~mapovým podkladom. Ako nadstavba bola plánovaná funkcionalita geofencingu -- definovanie virtuálnych zón s~automatickými upozorneniami pri vstupe alebo výstupe sledovaných osôb.

%% PREHĽAD NAJDÔLEŽITEJŠÍCH ZDROJOV %%

Pri práci na projekte som čerpal predovšetkým z~oficiálnych dokumentácií jednotlivých technológií. Pre backend vývoj to bola dokumentácia frameworku Flask~\cite{flask_docs} a~knižnice Socket.IO~\cite{socketio_docs}, pre databázovú vrstvu dokumentácia PostgreSQL~\cite{postgresql_docs}, TimescaleDB~\cite{timescaledb_docs} a~PostGIS~\cite{postgis_docs}. Na frontendovej strane som využíval dokumentáciu React~\cite{react_docs} a~mapovej knižnice Leaflet~\cite{leaflet_docs}. Kľúčovým zdrojom pre pochopenie IoT protokolov boli špecifikácie LoRaWAN~\cite{lorawan_spec} a~MQTT~\cite{mqtt_spec}, spolu s~datasheetmi trackerov SenseCap~\cite{sensecap_datasheet}. Nemenej dôležitá bola konzultácia s~mentorom Bc.~Janom Mitorajom, ktorý mi poskytoval praktické rady vychádzajúce z~jeho skúseností s~podobnými projektmi vo firme Railsformers.

%% ŠTRUKTÚRA PRÁCE %%

Práca je štruktúrovaná do šiestich hlavných kapitol. Prvá kapitola predstavuje firmu Railsformers~s.r.o. a~moje pracovné zaradenie v~rámci praxe. Druhá kapitola obsahuje prehľad zadaných úloh s~odhadom časovej náročnosti. Tretia kapitola, ktorá tvorí jadro práce, podrobne popisuje postup riešenia jednotlivých úloh od konfigurácie hardvéru cez implementáciu backendu a~frontendu až po testovanie. Štvrtá kapitola reflektuje teoretické a~praktické znalosti zo štúdia, ktoré som pri práci uplatnil. Piata kapitola naopak popisuje znalosti, ktoré mi chýbali a~ktoré som si musel počas praxe doplniť. Šiesta kapitola sumarizuje dosiahnuté výsledky a~obsahuje celkové zhodnotenie praxe. Záver potom prináša súhrnnú reflexiu a~návrhy na možné budúce rozšírenia systému.

\endinput
